\documentclass[11pt,a4paper,twoside]{report}

\usepackage[utf8]{inputenc}
\usepackage[T1]{fontenc}
\usepackage[ngerman]{babel}
\usepackage{amsmath}
\usepackage{amssymb}
\usepackage{geometry}
\geometry{left=2.5cm, right=2.5cm, top=2.5cm, bottom=2.5cm}
\usepackage{fancyhdr}
\usepackage{graphicx}
\usepackage{xcolor}
\usepackage{listings}
\usepackage{float}
\usepackage{algorithm}
\usepackage{algorithmic}
\usepackage{array}
\usepackage{booktabs}
\usepackage{hyperref}

% Code-Styling
\lstset{
    language=Python,
    basicstyle=\ttfamily\small,
    keywordstyle=\color{blue}\bfseries,
    commentstyle=\color{gray},
    stringstyle=\color{red},
    breaklines=true,
    showstringspaces=false,
    frame=single,
    backgroundcolor=\color{white},
    numberstyle=\tiny,
    numbers=left,
    stepnumber=1,
    columns=flexible
}

\pagestyle{fancy}
\fancyhf{}
\fancyhead[L]{Distributed Acoustic Ray Tracer}
\fancyhead[R]{\thepage}
\fancyfoot[C]{Ticket 3 \& Raumdaten-Integration}

\title{\Huge\textbf{Implementationsdokumentation}\\[0.3cm]
       \Large Audio I/O \& 3D-Raumdaten}
\author{Kais}
\date{\today}

\begin{document}

\maketitle
\tableofcontents
\newpage

% ========================================
\chapter{Einf\"uhrung}
% ========================================

Dieses Dokument beschreibt die Implementierung zweier kritischer Komponenten des Distributed Acoustic Ray Tracers:

\begin{enumerate}
    \item \textbf{Audio I/O (Ticket 3):} Einlesen trockener WAV-Dateien und Export des verarbeiteten, halligen Audio-Signals mit Clipping-Pr\"avention
    \item \textbf{Raumdaten (OBJ-Parser):} Laden von 3D-Modellen aus \texttt{.obj}-Dateien und Konvertierung in das interne JSON-Format
\end{enumerate}

Die Implementierung erfolgt in Python und integriert sich nahtlos in die bestehende Pipeline, die Rust-basierte Physik-Engine, DSP-Convolution und Distributed-Computing-Frameworks verbindet.

\section{Architektur-\"Ubersicht}

Das System folgt einer erweiterten Pipeline mit automatischer Raumgeometrie-Integration:

\begin{enumerate}
    \item[\textbf{[0/4]}] OBJ-Import: \texttt{room.obj} $\to$ \texttt{room\_geometry.json} mit Material-Absorption
    \item[\textbf{[1/4]}] Rust Physics Engine: L\"adt Raumgeometrie, simuliert Schallstrahlen, erzeugt IR
    \item[\textbf{[2/4]}] Audio-Einlesung: L\"adt trockene WAV-Datei
    \item[\textbf{[3/4]}] DSP-Convolution: Wendet IR mit frequenzabh\"angiger Material-D\"ampfung an
    \item[\textbf{[4/4]}] Audio-Export: Normalisiert und speichert Ergebnis
\end{enumerate}

Das Wandmaterial wird in \texttt{input\_config.json} konfiguriert und durchgehend durch alle Stufen propagiert: von der Geometrie-Absorption im Rust-Tracer bis zur oktavbandbasierten Frequenzformung in der Convolution.

% ========================================
\chapter{Audio I/O Handling}
% ========================================

\section{Anforderungen (Akzeptanzkriterien)}

\begin{enumerate}
    \item Sauberes Einlesen der \texttt{unprocessed.wav} ohne Qualit\"atsverlust
    \item Export der finalen WAV-Datei ohne Artefakte
    \item Normalisierung, damit geschichtetes Audio beim Export nicht clippt
\end{enumerate}

\section{Theoretischer Hintergrund}

\subsection{WAV-Format und Audio-Digitalisierung}

Das WAV-Format speichert Audiodaten als zeitdiskrete Samples mit konstanter Abtastfrequenz $f_s$ (typisch 44.1 kHz oder 48 kHz):

$$x[n] = \int_{(n-0.5)/f_s}^{(n+0.5)/f_s} x_{\text{analog}}(t) \, dt$$

Jedes Sample liegt im Bereich $[-1.0, 1.0]$ (bei normalisierten 16/32-Bit PCM).

\subsection{Clipping und Normalisierung}

Wenn w\"ahrend der Convolution (Faltung) die Amplitude $|y[n]|$ \"uber 1.0 steigt, tritt \textbf{Clipping} auf:

$$y_{\text{clipped}}[n] = \begin{cases}
    1.0 & \text{wenn } y[n] > 1.0 \\
    y[n] & \text{wenn } |y[n]| \leq 1.0 \\
    -1.0 & \text{wenn } y[n] < -1.0
\end{cases}$$

Um dies zu vermeiden, normalisieren wir:

$$y_{\text{norm}}[n] = \frac{y[n]}{A_{\max}} \quad \text{wobei} \quad A_{\max} = \max_n |y[n]|$$

\section{Implementierung: \texttt{wav\_handler.py}}

\subsection{WAV-Einlesung}

\begin{lstlisting}
import numpy as np
import soundfile as sf

def read_wav(wav_path: str) -> tuple:
    """
    Reads a WAV file and returns the audio data as a NumPy array
    along with the corresponding sample rate.
    """
    audio_data, sample_rate = sf.read(wav_path)
    return audio_data, sample_rate
\end{lstlisting}

\textbf{Funktionsweise:}

\begin{itemize}
    \item Die \texttt{soundfile}-Bibliothek liest WAV-Dateien robust
    \item \texttt{audio\_data} ist ein NumPy-Array der Form \texttt{(samples,)} oder \texttt{(samples, channels)}
    \item \texttt{sample\_rate} ist die Abtastfrequenz in Hz
\end{itemize}

\subsection{WAV-Export mit Normalisierung}

\begin{lstlisting}
def write_wav(wav_path: str, audio_data, sample_rate: int) -> None:
    """
    Normalizes the audio data to prevent clipping and exports it to a WAV file.
    """
    max_amplitude = np.max(np.abs(audio_data))
    if max_amplitude > 1.0:
        audio_data = audio_data / max_amplitude
    sf.write(wav_path, audio_data, sample_rate)
\end{lstlisting}

\textbf{Algorithmus:}

\begin{algorithm}
    \caption{Normalisierter WAV-Export}
    \begin{algorithmic}
        \STATE $A_{\max} \gets \max_n |y[n]|$
        \IF{$A_{\max} > 1.0$}
            \STATE $y[n] \gets y[n] / A_{\max}$ f\"ur alle $n$
        \ENDIF
        \STATE Schreibe normalisierte Samples in WAV-Datei
    \end{algorithmic}
\end{algorithm}

\subsection{Integration in die Pipeline}

\begin{lstlisting}
# main.py Schritt [2/4]: Audio laden
dry_audio, sample_rate = read_wav(input_wav_path)

# Schritt [3/4]: Convolution mit Material-Absorption
conv_material = ir_data.get("metadata", {}).get("wall_material")
wet_audio = convolve(dry_audio, sample_rate, ir_data, material=conv_material)

# Schritt [4/4]: Normalisieren und exportieren
write_wav(output_wav_path, wet_audio, sample_rate)
\end{lstlisting}

Der \texttt{wall\_material}-Wert wird automatisch aus \texttt{ir\_output.json} gelesen, das der Rust-Tracer nach der Simulation schreibt. Damit nutzen Tracer und Convolution stets dasselbe Material.

% ========================================
\chapter{3D-Raumdaten: OBJ-Parser}
% ========================================

\section{Theoretischer Hintergrund}

\subsection{OBJ-Format}

Das Wavefront \texttt{.obj}-Format ist ein textbasiertes 3D-Modellformat:

\begin{lstlisting}[language={}]
# Vertices (Ecken)
v  1.0  2.0 -1.5    # Vertex 1
v  2.0  2.0 -1.5    # Vertex 2
v  1.0  3.0 -1.5    # Vertex 3

# Faces (Dreiecke)
f  1/1/1  2/2/1  3/3/1
\end{lstlisting}

\textbf{Indexierung:} OBJ nutzt 1-basierte Indizes, Python 0-basierte -- Umrechnung n\"otig.

\subsection{Dreiecks-Repr\"asentation}

$$\text{Triangle} = \begin{bmatrix}
x_1 & y_1 & z_1 \\
x_2 & y_2 & z_2 \\
x_3 & y_3 & z_3
\end{bmatrix}$$

\section{Implementierung: \texttt{obj\_parser.py}}

\subsection{OBJ-Loader}

\begin{lstlisting}
def load_obj_to_triangles(file_path: str) -> list:
    vertices = []
    triangles = []
    path = Path(file_path)
    if not path.exists():
        raise FileNotFoundError(f"The 3D file '{file_path}' was not found.")
    with path.open('r') as f:
        for line in f:
            line = line.strip()
            if line.startswith('v '):
                parts = line.split()
                x, y, z = float(parts[1]), float(parts[2]), float(parts[3])
                vertices.append([x, y, z])
            elif line.startswith('f '):
                parts = line.split()
                v1_idx = int(parts[1].split('/')[0]) - 1
                v2_idx = int(parts[2].split('/')[0]) - 1
                v3_idx = int(parts[3].split('/')[0]) - 1
                triangle = [vertices[v1_idx], vertices[v2_idx], vertices[v3_idx]]
                triangles.append(triangle)
    return triangles
\end{lstlisting}

\subsection{Material-basierter Export: \texttt{export\_room\_geometry()}}

Die neue Exportfunktion verbindet OBJ-Geometrie mit der Materialdatenbank
(\texttt{materials.json}) und schreibt \texttt{room\_geometry.json} f\"ur den Rust-Tracer.
Jedes Dreieck erh\"alt die mittlere Absorptionskonstante des gew\"ahlten Materials:

\begin{lstlisting}
_MATERIALS_JSON = Path(__file__).resolve().parent.parent / "rust_service" / "materials.json"

def _mean_absorption(material: str) -> float:
    """Mittlere Absorption aus materials.json fuer das gewaehlte Material."""
    data = json.loads(_MATERIALS_JSON.read_text())
    alphas = data["materials"][material]["absorption"]
    return sum(alphas) / len(alphas)

def export_room_geometry(obj_path: str, material: str, out_path: str) -> None:
    triangles_raw = load_obj_to_triangles(obj_path)
    absorption = _mean_absorption(material)
    triangles = [
        {"v0": t[0], "v1": t[1], "v2": t[2], "absorption": absorption}
        for t in triangles_raw
    ]
    out = {"triangles": triangles, "triangle_count": len(triangles)}
    Path(out_path).write_text(json.dumps(out, indent=2))
\end{lstlisting}

\subsection{Ausgabeformat: \texttt{room\_geometry.json}}

Der Rust-Tracer liest Dreiecke inkl.\ Material-Absorption aus dieser Datei:

\begin{lstlisting}[language=JavaScript]
{
  "triangles": [
    {"v0": [0.0, 0.0, 0.0], "v1": [12.0, 0.0, 0.0],
     "v2": [12.0, 12.0, 0.0], "absorption": 0.033}
  ],
  "triangle_count": 12
}
\end{lstlisting}

Die \texttt{absorption}-Werte entsprechen dem arithmetischen Mittel aller 8 Oktavb\"ander
aus \texttt{materials.json}. Die frequenzabh\"angige Verarbeitung findet ausschlie\ss{}lich
in der Python-Convolution statt.

% ========================================
\chapter{Integration und Testing}
% ========================================

\section{End-to-End-Pipeline}

Die Pipeline verl\"auft nun in 5 Schritten. Schritt~0 ist vollautomatisch:
\texttt{room.obj} wird erkannt, mit dem konfigurierten Material in
\texttt{room\_geometry.json} umgewandelt und dem Rust-Tracer bereitgestellt.

\begin{lstlisting}
# Schritt 0: OBJ -> room_geometry.json (automatisch wenn room.obj vorhanden)
if Path("room.obj").exists():
    export_room_geometry("room.obj", wall_material,
                         "../rust_service/room_geometry.json")

# Schritt 1: Rust Physics Engine
ir_data = run_rust_and_load_json(rust_binary_path, json_path)

# Schritt 2: Audio laden
dry_audio, sample_rate = read_wav(input_wav_path)

# Schritt 3: Convolution mit Material-Absorption
conv_material = ir_data.get("metadata", {}).get("wall_material")
wet_audio = convolve(dry_audio, sample_rate, ir_data, material=conv_material)

# Schritt 4: Normalisieren und exportieren
write_wav(output_wav_path, wet_audio, sample_rate)
\end{lstlisting}

\section{Material-Konfiguration}

Das Wandmaterial wird in \texttt{rust\_service/input\_config.json} festgelegt:

\begin{lstlisting}[language=JavaScript]
{
  "wall_material": "concrete",
  "max_bounces": 100,
  "rays_to_cast": 1000
}
\end{lstlisting}

Vom Rust-Tracer wird \texttt{wall\_material} in \texttt{ir\_output.json} exportiert
und von der Python-Pipeline automatisch f\"ur die Convolution \"ubernommen.
Vegerf\"ugbare Materialien (aus \texttt{materials.json}):

\begin{table}[H]
    \centering
    \begin{tabular}{|l|l|c|}
        \hline
        \textbf{Material} & \textbf{Beschreibung} & \textbf{Mittl. Absorption} \\
        \hline
        \texttt{concrete} & Beton (stark reflektierend) & 0.033 \\
        \hline
        \texttt{drywall}  & Trockenbau (m\"a\ss{}ig) & 0.150 \\
        \hline
        \texttt{glass}    & Glas (stark reflektierend) & 0.026 \\
        \hline
        \texttt{carpet}   & Teppich (stark absorbierend) & 0.425 \\
        \hline
        \texttt{foam}     & Schaum (sehr stark absorbierend) & 0.650 \\
        \hline
    \end{tabular}
    \caption{Verf\"ugbare Materialien in \texttt{materials.json}}
\end{table}

\section{Abh\"angigkeiten}

\begin{table}[H]
    \centering
    \begin{tabular}{|l|l|c|}
        \hline
        \textbf{Modul} & \textbf{Bibliotheken} & \textbf{Version} \\
        \hline
        \texttt{wav\_handler.py} & soundfile, numpy & Any \\
        \hline
        \texttt{obj\_parser.py} & pathlib, json & Built-in \\
        \hline
        \texttt{convolution.py} & scipy, numpy & Any \\
        \hline
        \texttt{materials.json} & --- (Datenbasis) & --- \\
        \hline
    \end{tabular}
\end{table}

% ========================================
\chapter{Mathematische Modelle}
% ========================================

\section{Convolution im Frequenzraum}

$$Y(f) = X(f) \cdot H(f), \quad y[n] = \mathcal{F}^{-1}\{Y(f)\}$$

\section{Oktavband-basierte Material-D\"ampfung}

Der Rust-Tracer liefert pro Schallstrahl nur einen skalaren Druck. Die
frequenzabh\"angige D\"ampfung des Wandmaterials wird in der Convolution
nachtr\"aglich im Frequenzbereich modelliert.

Ein Tap zur Zeit $t$ hat circa $n(t) = t \,/\, t_{\text{mfp}}$ Reflexionen
erfahren. Pro Oktavband $b$ mit Absorptionskoeffizient $\alpha_b$ verbleibt
der Anteil:

$$\text{Gain}_b(t) = (1 - \alpha_b)^{n(t)} = (1 - \alpha_b)^{\,t / t_{\text{mfp}}}$$

Die Impulsantwort wird bandweise gefiltert und mit dem zeitvarianten Gain
gewichtet:

$$h_{\text{mat}}[n] = \sum_{b} \underbrace{(\text{Bandpass}_b * h)[n]}_{\text{gefilterte IR}} \cdot (1-\alpha_b)^{\,n/(f_s \cdot t_{\text{mfp}})}$$

Hohe B\"ander mit gro\ss{}em $\alpha_b$ (z.B.\ Teppich: $\alpha_{8\,\text{kHz}} = 0.70$)
klingen schnell aus --- der Raum klingt dunkler und trockener.
Bei Beton ($\alpha_b \approx 0.02\text{--}0.05$) bleibt der Klang lange hell.

\section{Ray-Tracer Intersection (M\"oller-Trumbore)}

Gegeben Strahl $\vec{r}(t) = \vec{o} + t\vec{d}$ und Dreieck $(\vec{a}, \vec{b}, \vec{c})$:

$$t = \frac{(\vec{c} - \vec{a}) \times (\vec{b} - \vec{a}) \cdot (\vec{a} - \vec{o})}{(\vec{c} - \vec{a}) \times (\vec{b} - \vec{a}) \cdot \vec{d}}, \quad \tau = \frac{t}{343 \text{ m/s}}, \quad n = \lfloor \tau \cdot f_s \rfloor$$

% ========================================
\chapter{Akzeptanzkriterien und Validierung}
% ========================================

\section{Ticket 3: Audio I/O}

\begin{table}[H]
    \centering
    \begin{tabular}{|p{4.5cm}|p{2cm}|p{5cm}|}
        \hline
        \textbf{Kriterium} & \textbf{Erf\"ullt?} & \textbf{Test} \\
        \hline
        Sauberes Einlesen .wav & \checkmark & \texttt{len(audio\_data) > 0} \\
        \hline
        Export ohne Qualit\"atsverlust & \checkmark & H\"ortest, Spektrum-Analyse \\
        \hline
        Normalisierung gegen Clipping & \checkmark & \texttt{max(abs(audio)) <= 1.0} \\
        \hline
    \end{tabular}
\end{table}

\section{Raumdaten: OBJ-Parser \& Material-Integration}

\begin{table}[H]
    \centering
    \begin{tabular}{|p{4.5cm}|p{2cm}|p{5cm}|}
        \hline
        \textbf{Kriterium} & \textbf{Erf\"ullt?} & \textbf{Test} \\
        \hline
        Parser liest .obj & \checkmark & \texttt{len(triangles) > 0} \\
        \hline
        JSON-Export mit Absorption & \checkmark & \texttt{absorption}-Feld pro Dreieck \\
        \hline
        Rust l\"adt OBJ-Geometrie & \checkmark & \texttt{room\_geometry.json} erkannt \\
        \hline
        Material durchgereicht & \checkmark & \texttt{wall\_material} in \texttt{ir\_output.json} \\
        \hline
        Frequenzabh\"ang. Convolution & \checkmark & \texttt{apply\_material\_absorption()} \\
        \hline
    \end{tabular}
\end{table}

% ========================================
\chapter{Zusammenfassung und Ausblick}
% ========================================

\section{Implementierte Features}

\begin{enumerate}
    \item \textbf{Audio I/O}: Robust WAV-Lesen/Schreiben mit soundfile
    \item \textbf{Clipping-Pr\"avention}: Peak-Normalisierung vor Export
    \item \textbf{OBJ-Parser}: Liest Vertices/Faces, konvertiert in Dreiecke
    \item \textbf{Material-Integration}: \texttt{export\_room\_geometry()} verbindet OBJ mit \texttt{materials.json}
    \item \textbf{Rust-OBJ-Pipeline}: Rust l\"adt \texttt{room\_geometry.json} automatisch
    \item \textbf{Material-Durchleitung}: \texttt{wall\_material} von Config bis Convolution
    \item \textbf{Frequenzabh\"angige D\"ampfung}: Oktavband-EQ via \texttt{apply\_material\_absorption()}
    \item \textbf{Pipeline-Integration}: Vollautomatisch ab Schritt~0 (OBJ-Erkennung)
\end{enumerate}

\section{M\"ogliche Erweiterungen}

\begin{enumerate}
    \item \textbf{Multi-Channel Audio}: Stereo/5.1 Surround Support
    \item \textbf{Per-Dreieck-Materialien}: Verschiedene Materialien pro Wandfl\"ache aus \texttt{.mtl}-Datei
    \item \textbf{Adaptive Normalisierung}: Loudness-basierte Pegel (LUFS) statt Peak
    \item \textbf{Geometrie-Optimierungen}: BVH-Tree f\"ur schnelle OBJ-Intersects
\end{enumerate}

\appendix
\chapter{Code-Referenz}

\section{wav\_handler.py}
Code-Referenz: \texttt{python\_service/wav\_handler.py}

\section{obj\_parser.py}
Code-Referenz: \texttt{python\_service/obj\_parser.py}

\section{main.py}
Code-Referenz: \texttt{python\_service/main.py}

\end{document}
